\scanpage{426}
% Main-text lines 1--3 continue the canceled draft on scan 425.
\deleted{also means \unsure{simply} that $\mathcal{X}^{\Gamma'}\ne\varnothing$ if $\mathcal{X}$ connected, or again that $\Gamma'$ \unsure{contains} a stabilizer $\Gamma_P$.} We assume

% Main-text line 4 and its small parenthetical addition.
\textcircled{\(1^\circ\)} $\forall P\in\mathcal{X}$, $\Gamma_P\ne\{1\}$, and $\Gamma_P$ of finite index \unsure{in $\Gamma$} \inserted{(i.e. the orbit of $P$ finite)}

% Main-text lines 5--21.
We want conditions which ensure that $\mathcal{X}$ and the $\Pi_{U_I}$ \deleted{\unsure{for}} etc, \unsure{can|may} be reconstructed up to \deleted{\unsure{the help}} canonical \unsure{isomorphism} with the help of the purely groupoid-theoretic ``bare'' or topossic \inserted{(with operators $\Gamma$)} \unsure{corresponding|corresponding} data. \deleted{We \illegible{1}} We can say it in learned topossic language, but we are going to express it in terms of the theory of groups with loops. Let $I$ be a subset of $\mathcal{X}$ stable under \unsure{$\Gamma^*$|$\Gamma'$}, we would like to lay down conditions which \deleted{\unsure{we take} a \unsure{subgroup}} \unsure{allow|will allow} us to express \inserted{\unsure{(for such $I$)}} the situation of the groupoids \inserted{$\Pi_V$} associated with the $V\supset U_I$ (i.e. the $U_J$ with $J\subset I$), \inserted{the $\Pi_{D_i^*}$ $(i\in I)$,} and the operation of $\Gamma$ on them, in terms of the \deleted{\unsure{operation}} sole operation of $\Gamma$ on the outer group with loops \deleted{$\pi_I$ associated} $\pi_1(\Pi_{U_I})$ associated with $I$ ---
\scanpage{427}
\mspage{216}
\margin{\unsure{$\rightsquigarrow$|$\longrightarrow$}}
% Main-text lines 1--5.
or rather (for this is immediate, by the operation of drilling holes), \inserted{in the reverse direction,} show how,
\deleted{\unsure{for a} $J\subset I$ \unsure{invariant} under $\Gamma$}
\deleted{\illegible{5}}
\deleted{\unsure{reconstruct} \illegible{2},}
% Main-text lines 6--11.
subject to anabelianity \unsure{$\zeta$|$\xi$} say of \unsure{$\Pi_{\mathcal X}$|$B_{\Pi_{\mathcal X}}$} (to fix ideas), we can $\pm$ \unsure{reconstruct} $\Pi_{U_I}$ and the action of $\Gamma$ on it \unsure{....|...},
starting from \unsure{$\Pi_{\mathcal X}$} \inserted{(or rather, \deleted{from}\inserted{from the \unsure{topos}} \unsure{$B_{\Pi_{\mathcal X}}$})} and from the action of $\Gamma$ on it.
We want at least an \inserted{\unsure{intrinsic}} description of the \unsure{elements} of \unsure{$\mathcal X$}, \inserted{and of the action of $\Gamma$ on it,} via
% Main-text lines 12--15.
this ``soft'' action --- which, for $X$ connected, again amounts to an outer action of $\Gamma$ on a group with loops (\uline{without} loops!) $\pi$...

% Main-text lines 16--21; the parenthesis continues on the next scan.
A first condition, in weak form, \unsure{it is|is} that (assuming \unsure{$\mathcal X$} connected, and calling : $\pi(I)$ the outer group $\pi_1(\Pi_{U_I})$, for \unsure{every} finite subset $I$ of \unsure{$\mathcal X$}) that \unsure{for} $\pi(I)$ anabelian (which will be the case for $I$ ``large enough'',
\scanpage{428}
% Main-text lines 1--3.
at least if \unsure{$\mathcal X$} is infinite \unsure{so|of} that we can take $I$ of arbitrarily large cardinality...) \deleted{we have} the obvious map
% Deleted display and replacement display, approximately lines 4--5.
\[
\deleted{\unsure{\mathrm{Pt}(B_{\pi(I),\Gamma\natural})|\mathrm{Pt}(B_{\pi(I)},\Gamma\natural)}\longrightarrow\text{classes of}}
\]
\[
2^\circ)\quad
\deleted{\unsure{\mathcal X}}U_I\longrightarrow
\unsure{\mathrm{Pt}(B_{\pi(I),\Gamma\natural})|\mathrm{Pt}(B_{\pi(I)},\Gamma\natural)}
\simeq
\begin{gathered}
\unsure{\inserted{\text{classes of }\pi(I)\text{-conjugation }\unsure{\text{of}|\text{of the}}}}\\
\unsure{\text{germs}}\text{ of}\\
\deleted{\text{of}}\text{ splittings of}\\
\text{the extension }\mathcal E(I)\\
\text{of }\deleted{\unsure{\mathcal E}}\Gamma\text{ by }\pi(I)
\end{gathered}
\]
% Main-text lines 6--15.
\unsure{(compatible|compatible} with the actions of $\Gamma$) be \uline{injective}. But to go further, we would have to be able to give a characterization of the image. Let us note (placing ourselves henceforth in the profinite context) that if we give ourselves on the \unsure{$\Pi(I)$|$\pi(I)$} some \unsure{arithmetizations}, \unsure{compatible} with the outer homomorphisms of plugging holes $\pi(I)\longrightarrow\pi(J)$, and with the action of $\Gamma$ on the $\pi(I)$, we therefore have a
% Main-text lines 16--20.
canonical homomorphism from $\Gamma$ into the common group \deleted{of the} $\boldsymbol{\Gamma}$ of \unsure{arithmetically} \unsure{(}outer automorphisms of these \unsure{arithmetizations}.
(and even $\Gamma\longrightarrow$ \unsure{$M_{\pi(\varnothing)=\pi}$}, --- which is \inserted{already} a markedly stronger \inserted{datum}). We can therefore
\scanpage{429}
\mspage{217}
% Main text, approximate lines 1--5.
suppose \unsure{given} a notion of \uuline{admissible}
\uline{lifting} of such an outer arithmetic
action, in such a way that the
preceding \unsure{morphism} is a bijection
\[
 U_I=\unsure{\mathcal X|\mathfrak X}\setminus I
 \xrightarrow{\sim}
 \mathrm{Pt}_{\mathrm{ad}}\bigl(B_{\pi(I)},\Gamma\natural\bigr)
\]
This means moreover, in practice,
that the notion of admissibility satisfies
conditions $1^\circ)$ : $6^\circ)$ \unsure{seen} above.

% Main text, approximate lines 9--12.
As for knowing \unsure{what} should be
called \unsure{operations} \deleted{\unsure{outer}} ``\unsure{admissible|admissible}'' of
$\Gamma$ on an \inserted{outer} \unsure{(}profinite\unsure{)|,} group with loops,
this remains conjectural for the moment.

% Main text, approximate lines 13--21; interlinear revisions retained.
We could by way of experiment
\unsure{conjecture} \deleted{that \unsure{every} \illegible{2}}
\inserted{that the following notion would work}
\unsure{Let us call|Let us call}
``admissible'' every such action
\[
 \Gamma\longrightarrow\widehat{\widehat{\mathcal T}}(\pi)
\]
which respects an arithmetization $a$,
i.e. which factors through the corresponding $N_a$, \unsure{and} such that the image of
$\Gamma\longrightarrow N_a/\widehat{\mathcal T}_a=\boldsymbol{\Gamma}_a$
is \unsure{open|open},
\deleted{\inserted{and that $\pi^{\Gamma'}=(1)$}}
\deleted{\unsure{the conditions}}, \inserted{and that}
\scanpage{430}
% Main text, approximate lines 1--5.
\deleted{\unsure{considered} above \unsure{$1^\circ:6^\circ$} \unsure{are} \unsure{satisfied|satisfied}}

$\forall$ \unsure{homomorphism} of plugging holes at $i\in I(\pi)$,
$\pi\longrightarrow\pi_i'$, $N(\pi)_i\longrightarrow M(\pi_i')$, the \unsure{corresponding} homomorphism
$\Gamma_i$, nor of any open subgroup $\Gamma'$
of $\Gamma_i$, does \unsure{not} normalize an $L_j'$ $(j\in I(\pi_i')=I\setminus\{i\})$.

% Upper left margin; several long strokes run through the whole note.
\margin{\deleted{
this condition is
perhaps \illegible{2} and \unsure{premature}
--- it could be
that for \unsure{every pair|all pairs}
of distinct indices $i,j$ of $I=I(\pi)$,
\inserted{and for the \unsure{corresponding} $L_i,L_j$},
the subgroup \deleted{\illegible{1}} \unsure{$K$} of
$M(\pi)$ \unsure{formed} of the
$g\in M(\pi)$ which \uline{normalize}
$L_i$, and which moreover normalize
$L_j$ \uline{mod} $\mathrm{Ker}(\unsure{\pi\longrightarrow\pi_i'})$,
is ``\unsure{too small}'' \deleted{\unsure{for}} \inserted{\unsure{i.e.}}
that its image in
$\unsure{\boldsymbol{\Gamma}|\Pi}=M(\pi)/\unsure{\widehat{\mathcal T}(\pi)|\widehat{\mathcal G}(\pi)}$
\uline{\unsure{is not}}
of finite index
}}

% Main text, approximate lines 6--14.
\unsure{perhaps} we would even have to impose
$\pi^{\unsure{(\Gamma^{\prime+})|\langle\Gamma^{\prime+}\rangle}}=(1)$, where $\Gamma^{\prime+}$ is the kernel of
$\Gamma'\xrightarrow{\chi}\widehat{\mathbb Z}^{*}$ --- in any case we expect
\unsure{---} that we then have $\pi^{\Gamma'}=(1)$,
and if this were not the case, we would have to
introduce it into the definition --- for
each operation of plugging holes
\unsure{relative} to $I'\subset I(\pi)$, and to the choice of an
$i\in I'$, allowing us to define $\Gamma_{(I',i)}\longrightarrow M(\pi(I'))$.

% Lower left marginal note.
\margin{
More generally,
for $i\in I'\subset I=I(\pi)$,
\deleted{\unsure{it is necessary}} \inserted{\unsure{we impose}} that the
true action of
\unsure{$\Gamma_{I',i}$} on \unsure{$\pi_i(I')$|$\Pi_i(I')$}
\inserted{\unsure{corresponding}} \unsure{to $i$}
normalize no
subgroup $L_j$
$(j\in I\setminus I')$ --- we
are reduced for
this to the case where
$I'=I\setminus\{j\}$, hence
where \unsure{$\pi(I')$|$\Pi(I')$} has \unsure{(only a} single \unsure{class} loops\dots\
(\unsure{if|and} $g\geq1$)
}

% Main text, approximate lines 15--21.
Let us note that a true action of $\Gamma$ on $\pi$
(not \unsure{only|only} \unsure{outer}) which lifts a
given\unsure{$+$} admissible action, is \inserted{\unsure{it too}} admissible
(in the sense that it defines an
\deleted{outer} \uline{admissible} action on a $\pi'$ of
type $(g,\nu+1)$), iff this action \inserted{\unsure{(rather its germ)}} normalizes
no $L_i$ $(i\in I(\pi))$, and that